\section{Exact Potts filters and independent component tensors}
\label{sec:potts-integration}

Report 17 contributes an optional scalar obstruction before the complete
homology fallback. The integrated selectors are \code{potts5},
\code{potts-exact}, and \code{potts-exact-factorized}; the existing
\code{matching} Jones backend remains the default. These selectors are
independent of the Khovanov backend and of adaptive cancellation.
Potts representations of Jones evaluations and width-sensitive contraction
are established techniques, including the tensor implementation of
Meichanetzidis--Kourtis \cite{potts-tensor} and Maria's parameterized analysis
of quantum invariants \cite{quantum-width}. The contribution here is the
audited implementation, exact coefficient representation, component factoring,
and explicit resource and fallback contracts.

\subsection{Normalization without approximate algebraic numbers}
Let $D$ be a nonempty validated plane knot diagram, $G$ either checkerboard
Tait multigraph, $v$ its number of vertices, and $r$ the diagram writhe.
Loops and parallel edges are retained. At crossing $e$, let
$\epsilon_e\in\{1,-1\}$ be the bracket exponent of the smoothing that
omits its Tait edge; put $s=\sum_e\epsilon_e$.
A subset $S$ of included edges produces
\[
 c(S)=2k(S)-v+|S|
\]
smoothing circles. Indeed a regular neighborhood of the spanning plane
subgraph has Euler characteristic $v-|S|=2k(S)-c(S)$, including its isolated
vertices. With $\delta=-A^2-A^{-2}$ and the full bracket of one circle equal
to $\delta$, expansion of the smoothing sum gives
\[
 B_D(A)=\delta^{-v}A^s
 \sum_{S\subseteq E}(\delta^2)^{k(S)}
       \prod_{e\in S}\delta A^{-2\epsilon_e}.
\]
At the integer specialization $q=\delta^2$, expanding a spin product shows
that the sum is a partition function: each selected edge requires equal
endpoint colors, so the assignments constant on $(V,S)$ number $q^{k(S)}$.
Writing $x=A^4$ gives
\[
 q=x+2+x^{-1},\qquad
 Z_G=\sum_{\sigma:V\to[q]}\prod_{e=uv}
 \begin{cases}-x^{-\epsilon_e},&\sigma(u)=\sigma(v),\\
               1,&\sigma(u)\ne\sigma(v).
 \end{cases}
\]

The normalized Jones polynomial is $V_D(A^{-4})=B_D(A)/
\bigl(\delta(-A^3)^r\bigr)$. Substitution of
$\delta=-(x+1)A^{-2}$ yields
\begin{equation}\label{eq:potts-normalization}
 V_D(x^{-1})=(-1)^{r+v+1}x^h\frac{Z_G}{(x+1)^{v+1}},
 \qquad h=\frac{s-3r+2(v+1)}4.
\end{equation}
For a knot, $h$ is an integer. The residue modulo four of
$s_\sigma-3r+2(c_\sigma-1)$ is unchanged by a smoothing flip.
In the Seifert state $s_\sigma=r$ and $c_\sigma-1\equiv r\pmod2$;
in the all-omitted state $c_\sigma=v$, giving the assertion.
Thus the exact obstruction compares $Z_G$ with
\[
 Z_U=(-1)^{r+v+1}x^{-h}(x+1)^{v+1}.
\]
No fourth root or division by $x+1$ is computed. The crossing-free knot
is handled separately with $Z_G=Z_U=q$.

For every integer $q\ge5$, work in
$R_q=\mathbb Z[x]/(x^2-(q-2)x+1)$.
The discriminant lies strictly between $(q-3)^2$ and $(q-2)^2$, so this
polynomial is irreducible and each element has a unique integer pair
$(a,b)$ representing $a+bx$. Here $x^{-1}=q-2-x$ and
\[
 (a,b)(c,d)=(ac-bd,\ ad+bc+(q-2)bd).
\]
Multiplication by an equal-spin edge uses just additions and multiplication
by $q-2$:
\[
 -x^{-1}(a,b)=(-(q-2)a-b,a),\qquad
 -x(a,b)=(b,-a-(q-2)b).
\]
Partition inequality therefore certifies a nontrivial Jones value exactly.
Partition equality is inconclusive. The exact implementations default to
$q=6$, whose roots are $2\pm\sqrt3$; no numerical embedding is chosen.
Witness pairs are serialized as signed hexadecimal strings, avoiding
Python's limit on conversion of very large integers to decimal strings.

\subsection{A half-frontier bound for every supplied crossing order}
At a cut in the crossing order, let $w_i$ count shadow edges joining a
processed crossing to an unprocessed one. A Tait vertex is active when it
is incident to both processed and unprocessed edges; write $f_i$ for their
number. Traverse the cyclic boundary walk of an active selected-color face.
The processed/unprocessed marks change at least twice. Each change traverses
a cut shadow edge. A cut edge has just one side of the selected color and
can be charged to just that face. Hence
\begin{equation}\label{eq:potts-half-frontier}
 2f_i\le w_i.
\end{equation}
This argument allows repeated vertices and edges in face walks, and does
not require a common-disk ordering or a Catalan matching hypothesis.

The single-table algorithm introduces a spin at its first incident edge,
applies the edge weight, and forgets spins at their last incident edge.
Its key records the equality pattern of active colors, canonically naming
classes in first-occurrence order. Coefficients aggregate \emph{all} labeled
assignments in an orbit. An introduced spin can join any existing class
with multiplicity one, or create a new class with multiplicity $q-k$ when
$k$ classes already occur. Restriction and recanonicalization implement
forgetting. There are at most $q^{f_i}$ keys; an edge introduces at most
two extra vertices. At fixed $q$, the work is
\[
 \operatorname{poly}(n,w)\,q^{w/2},\qquad w=\max_i w_i,
\]
including polynomial bit arithmetic. To justify the latter, each coefficient
after $i$ edges and $I$ introduced vertices is a sum of at most $q^I$
assignments with pair norm bounded by $(q-1)^i$. Therefore its bit length
is $O(n\log q)$, since $v\le n+1$. The normalization pair also has that
bit bound. This is a scalar-query bound, not a bound on full homology.

\subsection{Independent processed components and exact orbit division}
The factored implementation keeps one tensor for each connected component
of processed Tait edges. A vertex forgotten from the frontier does not
erase interactions through it. Consequently this decomposition follows
processed connectivity, not connectivity of just the active vertices.
Relative color equalities between two components become relevant only when
an edge joins them. A closed component contributes a scalar factor.

For a pattern using $k$ colors, write $(q)_k=q(q-1)\cdots(q-k+1)$.
Its aggregate coefficient $C$ is $(q)_k$ times its per-assignment value,
because simultaneous color permutation preserves the partial partition sum.
When joining patterns with $k$ and $\ell$ classes, enumerate partial
bijections between their classes. An alignment with $K\le q$ total classes
has pre-edge aggregate coefficient
\begin{equation}\label{eq:potts-orbit-merge}
 \frac{C_L}{(q)_k}\frac{C_R}{(q)_\ell}(q)_K.
\end{equation}
The integer coordinates of each division must have zero remainder; the code
checks this before multiplying pairs. The operation is exact integer
division, not division in a finite field or a numerical approximation.
Each combined equality pattern determines its two restrictions and alignment
uniquely, proving that the merge neither loses nor duplicates assignments.

There are at most $q!$ relative alignments per pair of keys. Let $g$ be the
largest number of active vertices in any one processed component after an
edge. Only the current edge's endpoints can be forgotten at that step,
so a merged component has at most $g+2$ vertices before forgetting. Thus
the Cartesian product of its two input key sets is bounded by $q^{g+2}$,
not merely by the weaker $q^{2g}$. The resulting bound is
\[
 \operatorname{poly}(n,g)\,q!\,q^{g+2}
\]
coefficient and key work, with $O(nq^g)$ logical keys and polynomial
coefficient lengths. Here $g\le\max_i f_i\le w/2$.
The extra alignment work can dominate when factoring saves few states.

\subsection{A specialization collision that changing primes cannot repair}
Report 17 analyzes $W_m=\widehat{(\sigma_1\sigma_2^{-1})^m}$.
It is a knot exactly when $3\nmid m$. Its three-braid Jones trace identity
gives, at $t+t^{-1}=q-2$,
\[
 V_{W_m}(t)=q-2+T_m,\quad
 T_0=2,\quad T_1=3-q,\quad T_{m+1}=(3-q)T_m-T_{m-1}.
\]
At $q=5$ this is $3+2(-1)^m$, equal to one for every odd $m$.
In particular the nontrivial members with odd $m\ge5$ are missed even in
characteristic zero. Changing primes or using the conjugate root preserves
this equality. For $q\ge6$, set $a=q-3\ge3$ and
$U_m=(-1)^mT_m$. Then $U_0=2$, $U_1=a$,
$U_2=a^2-2>a$, and the recurrence proves strict increase thereafter.
Thus $a+1+(-1)^mU_m\ne1$ for every $m\ge2$.
For example, $V_{W_5}=-119$ at $q=6$.
This proves detection for this family only. The existing short-braid stage
already decides these inputs when their source braid is available.

\subsection{Integration, budgets, and validation}
The optional filter runs in the same place as the existing matching Jones
filter and reuses its crossing order. A mismatch returns \code{KNOTTED};
equality or local state/transition exhaustion continues the existing
Alexander, simplification, window, and complete homology stages.
The global cooperative deadline still returns \code{UNKNOWN}.
\code{--jones-backend} selects the filter in recognition, while
\code{jones --backend} selects a standalone query.
The standalone query now exits 3 with an inconclusive result when its
resource limit is reached: a nonempty error dictionary is never interpreted
as an obstruction. Invalid options exit 2.

State limits count represented keys. Inputs and a replacement table can
coexist during a merge, and normalized lists and alignment caches consume
additional memory, so this is not a resident-memory limit. Transition
limits count actual orbit updates, including relative alignments.
At fixed $q\le6$ the alignment cache has finitely many bounded entries;
larger color counts use lazy alignment generation.
The factored wrapper additionally checks the global deadline after final
normalization. The archived source remains unchanged.

The archived suite passed 172 tests; the first integrated checkpoint passed
315 tests. The retained independent full-smoothing oracle constructs Laurent
Jones polynomials without Tait graphs or production ring multiplication.
It checked 248 diagrams and 30,116 cube states, followed by 4,464 comparisons
per exact backend over three color counts, two shadings, and three orders.
Another 45 weaving checks reached 98 crossings. These are correctness
observations, not complexity proofs. The report's certificate replayer,
which recomputes the same backend, is not presented as an independent oracle.
The ring-valued Fitting proposals in report 17 are not integrated by this change.

\subsection{Measurements and the quasi-polynomial objective}
\label{sec:potts-measurements}
The reproducible driver \path{fast/benchmark_potts.py} measures seven
shuffled paired rounds with an exact-$q=6$ A/A control. Kernel timings
exclude construction of their common supplied order; the fragmented-order
case is deliberately adverse. A separate normal-pipeline comparison includes
diagram construction, ordering, cheap certificates, and fallbacks. Matching
and modular $q=5$ arms evaluate different specializations, so their timing
ratios alone do not establish equivalent detection power.
Detailed measurements are recorded in
\path{fast/results/potts_20261008.json}.

\begin{center}\small
\begin{tabular}{lrrrr}
\toprule
Input & Exact keys & Factored keys & Exact/factored & Pipeline gain\\
\midrule
Weaving, 10 crossings & 5 & 6 & 0.799 & bypassed\\
Weaving, 98 crossings & 5 & 6 & 0.934 & bypassed\\
Fragmented order, 16 crossings & 4111 & 205 & 23.612 & bypassed\\
Conway & 5 & 6 & 0.803 & 1.143\\
Kinoshita--Terasaka & 5 & 6 & 0.734 & 1.152\\
Stress braid, 36 crossings & 15 & 16 & 0.787 & bypassed\\
Hard unknot, 8 crossings & 2 & 3 & 0.749 & bypassed\\
\bottomrule
\end{tabular}
\end{center}
Here ``Exact/factored'' is the median paired kernel-time ratio; a value
above one favors factoring. ``Pipeline gain'' compares default matching
recognition with the single-table exact-$q=6$ option, including construction
and all earlier stages. On Conway and Kinoshita--Terasaka, every compared
pipeline reached its Jones filter and returned \code{KNOTTED}. The factored
option's corresponding pipeline ratios were only 1.069 and 1.053.
The other five cases were decided earlier, so their pipeline timings measure
noise and common stages rather than a benefit from Potts evaluation.

Factoring reduced transitions from 30,197 to 1,558 on the fixed fragmented
order but was slower in all six other kernel cases. The exact A/A ratios
ranged from 0.973 to 1.159, making small timing changes unsuitable for a
broad speedup claim. The robust large saving on the adverse order supports
keeping the factored implementation available; these data do not support
making it the default. Nor does the modest normal-pipeline improvement on
two fixtures establish superiority of the six-color specialization on
unseen diagrams.

At fixed $q$, a supplied order with $w=O(\log^2 n)$ makes this filter
quasi-polynomial. The stronger component parameter $g$ can make a poor
global frontier affordable. Neither guarantees a suitable order for every
diagram, nor that the scalar differs from the unknot value. The fallback
therefore retains its exponential worst case. This integration adds options
for a future adaptive scalar-query policy; it does not change the existing
progress-based adaptive Khovanov cancellation policy or promote an unmeasured
automatic Potts switch.
